%% kapa: exam=1 round=23 year=2012 period=1 form=A subject=civil_law subject_ko=민법 seq=3 number=3
%% source: Q-Net 감정평가사 기출문제 (cst00309) artlSeq=5052578 file=raw/1008/5052578/2012년 제23회 감정평가사 제1차시험 기출문제(1교시 A형).pdf
%% conversion: kordoc 4.24.51 -> tools/gen1.py (auto); review=proofread-sonnet; flags=none
\begin{question}{3}
조건과 기한에 관한 설명으로 옳지 않은 것은? (다툼이 있으면 판례에 의함)
\begin{choices}
  \item 기성조건이 정지조건이면 조건 없는 법률행위가 되고, 해제조건이면 그 법률행위는 무효이다.
  \item 조건의 성취로 인하여 불이익을 받을 당사자가 신의성실에 반하여 조건의 성취를 방해한 경우, 조건이 성취된 것으로 의제되는 시점은 방해행위 시점이다.
  \item 기한이익의 상실약정이 있는 경우, 특별한 사정이 없는 한 채권자가 통지나 청구 등을 하여야만 채무자의 기한이익이 상실된다.
  \item 부첩관계 종료를 해제조건으로 하여 부동산을 증여한 경우, 증여계약은 무효이다.
  \item 조건성취의 효력은 조건이 성취한 때로부터 생기고 소급하지 않는 것이 원칙이다.
\end{choices}
\end{question}
\answer{2}
